\paragraph{局部化有限商、周期动力学、容量原子恢复与缺陷重整化}
在 prime-register 的有限商分类、局部化 solenoid 的周期谱、bin-fold 临界容量的原子 defect law 以及 fold-gauge 常数列的前两阶 Bernoulli 模态已经建立之后，还可进一步抽取出下列六条彼此闭合的结构后果。

\begin{theorem}[局部化数系有限商的精确外部化公式]\label{thm:xi-localized-quotient-exact-sfree-externalization}
\leanverified{paper\_xi\_localized\_quotient\_exact\_sfree\_externalization}
设 \(S\subset\mathbb P\) 为有限素数集，并记
\[
G_S:=\ZZ[S^{-1}].
\]
对任意 \(n\ge 1\)，令
\[
n_S:=\prod_{p\in S}p^{v_p(n)},
\qquad
n_{S^\perp}:=\frac{n}{n_S}.
\]
则有自然同构
\[
\boxed{\;
G_S/nG_S\cong \ZZ/n_{S^\perp}\ZZ.\;}
\]
特别地，\(G_S\) 的每一个有限商都只由 \(n\) 的 \(S\)-自由部分决定。
\end{theorem}
\begin{proof}
对每个 \(p\in S\)，元素 \(p^{-1}\in G_S\) 存在，故乘法映射
\[
x\longmapsto px
\]
在 \(G_S\) 上可逆。于是对
\[
n=n_S\,n_{S^\perp}
\]
有
\[
nG_S=n_Sn_{S^\perp}G_S=n_{S^\perp}G_S.
\]
因此只需证明当
\[
\gcd\!\Bigl(n,\prod_{p\in S}p\Bigr)=1
\]
时，
\[
G_S/nG_S\cong \ZZ/n\ZZ.
\]

任取 \(x\in G_S\)，可写成
\[
x=\frac{a}{s},
\qquad
a\in\ZZ,
\]
其中 \(s\) 的全部素因子都属于 \(S\)。由于 \(\gcd(s,n)=1\)，\(s\) 在 \(\ZZ/n\ZZ\) 中可逆。定义
\[
\pi_n:G_S\longrightarrow \ZZ/n\ZZ,
\qquad
\pi_n\!\left(\frac{a}{s}\right):=a\,s^{-1}\pmod n.
\]
若 \(a/s=a'/s'\) 于 \(G_S\) 中相等，则 \(as'=a's\)，从而
\[
a\,s^{-1}\equiv a'\,s'^{-1}\pmod n,
\]
故 \(\pi_n\) 良定义。又因整数子群 \(\ZZ\subset G_S\) 已映满 \(\ZZ/n\ZZ\)，故 \(\pi_n\) 为满射。

再看其核。若
\[
\pi_n\!\left(\frac{a}{s}\right)=0,
\]
则 \(n\mid a\)。写 \(a=nb\)，便有
\[
\frac{a}{s}=n\cdot \frac{b}{s}\in nG_S.
\]
故
\[
\ker\pi_n\subset nG_S.
\]
反向包含显然成立，于是
\[
\ker\pi_n=nG_S.
\]
由第一同构定理得到
\[
G_S/nG_S\cong \ZZ/n\ZZ.
\]
最后把 \(n\) 替换为 \(n_{S^\perp}\) 即得结论。证毕。
\end{proof}

\begin{theorem}[局部化 solenoid 自覆盖的固定点公式与 Artin--Mazur \texorpdfstring{$\zeta$}{zeta} 函数]\label{thm:xi-localized-solenoid-fixedpoints-artin-mazur-series}
\leanverified{paper\_xi\_localized\_solenoid\_fixedpoints\_artin\_mazur\_series}
沿用定理 \ref{thm:xi-localized-quotient-exact-sfree-externalization} 的记号，并记
\[
\Sigma_S:=\widehat{G_S}.
\]
设整数 \(a\ge 2\) 满足
\[
\gcd\!\Bigl(a,\prod_{p\in S}p\Bigr)=1,
\]
令
\[
f_a:\Sigma_S\longrightarrow \Sigma_S
\]
为离散端“乘以 \(a\)”的 Pontryagin 对偶。则对每个 \(k\ge 1\)，有
\[
\boxed{\;
\mathrm{Fix}(f_a^k)\cong \widehat{G_S/(a^k-1)G_S},
\qquad
\#\mathrm{Fix}(f_a^k)=(a^k-1)_{S^\perp}.\;}
\]
并且其 Artin--Mazur \(\zeta\)-函数满足
\[
\boxed{\;
\zeta_{f_a}(z)
:=
\exp\!\left(
\sum_{k\ge 1}\frac{\#\mathrm{Fix}(f_a^k)}{k}z^k
\right)
=
\exp\!\left(
\sum_{k\ge 1}\frac{(a^k-1)_{S^\perp}}{k}z^k
\right).\;}
\]
特别地，当 \(S=\varnothing\) 时，
\[
\boxed{\;
\zeta_{f_a}(z)=\frac{1-z}{1-az}.\;}
\]
\end{theorem}
\begin{proof}
由对偶定义，对任意 \(x\in \Sigma_S\) 与 \(g\in G_S\)，有
\[
(f_a^k x)(g)=x(a^k g).
\]
因此
\[
x\in \mathrm{Fix}(f_a^k)
\iff
x(a^k g)=x(g)\quad(\forall g\in G_S),
\]
亦即
\[
x\bigl((a^k-1)g\bigr)=1\quad(\forall g\in G_S).
\]
这等价于 \(x\) 因子化经过商群
\[
G_S/(a^k-1)G_S.
\]
故有自然同构
\[
\mathrm{Fix}(f_a^k)\cong \widehat{G_S/(a^k-1)G_S}.
\]

再由定理 \ref{thm:xi-localized-quotient-exact-sfree-externalization}，
\[
G_S/(a^k-1)G_S\cong \ZZ/(a^k-1)_{S^\perp}\ZZ.
\]
有限循环群的 Pontryagin 对偶仍为同阶循环群，于是
\[
\#\mathrm{Fix}(f_a^k)=(a^k-1)_{S^\perp}.
\]
把该计数代入 Artin--Mazur \(\zeta\)-函数的定义，即得所示显式级数。

当 \(S=\varnothing\) 时，\((a^k-1)_{S^\perp}=a^k-1\)，故
\[
\log\zeta_{f_a}(z)
=
\sum_{k\ge 1}\frac{a^k-1}{k}z^k
=
-\log(1-az)+\log(1-z).
\]
指数化即得
\[
\zeta_{f_a}(z)=\frac{1-z}{1-az}.
\]
证毕。
\end{proof}

\begin{theorem}[bin-fold 极限容量曲线的原子测度唯一恢复]\label{thm:xi-binfold-capacity-atomic-measure-unique-recovery}
把定理 \ref{thm:xi-binfold-critical-capacity-defect-atomic-law} 的极限容量曲线写为
\[
\widehat C_\infty(s)
:=
\frac{\varphi}{\sqrt5}\min\{1,s\}
+
\frac{1}{\sqrt5}\min\{\varphi^{-1},s\}
\qquad (s>0).
\]
则存在唯一有限正测度
\[
\boxed{\;
\eta_\infty
:=
\frac{1}{\sqrt5}\delta_{\varphi^{-1}}
+
\frac{\varphi}{\sqrt5}\delta_{1}\;}
\]
使得对一切 \(s>0\) 都有
\[
\boxed{\;
\widehat C_\infty(s)=\int_{(0,\infty)}\min\{s,y\}\,\eta_\infty(\dd y).\;}
\]
并且在分布意义下，
\[
\boxed{\;
\frac{\dd}{\dd s}\widehat C_\infty(s)=\eta_\infty([s,\infty)),
\qquad
-\frac{\dd^2}{\dd s^2}\widehat C_\infty=\eta_\infty.\;}
\]
\end{theorem}
\begin{proof}
取
\[
\eta_\infty
:=
\frac{1}{\sqrt5}\delta_{\varphi^{-1}}
+
\frac{\varphi}{\sqrt5}\delta_1.
\]
则直接计算得到
\[
\int_{(0,\infty)}\min\{s,y\}\,\eta_\infty(\dd y)
=
\frac{1}{\sqrt5}\min\{s,\varphi^{-1}\}
+
\frac{\varphi}{\sqrt5}\min\{s,1\}
=
\widehat C_\infty(s).
\]
故积分表示成立。

再对各分段求导，得到
\[
\widehat C_\infty'(s)=
\begin{cases}
\dfrac{\varphi+1}{\sqrt5},& 0<s<\varphi^{-1},\\[2mm]
\dfrac{\varphi}{\sqrt5},& \varphi^{-1}<s<1,\\[2mm]
0,& s>1.
\end{cases}
\]
另一方面，
\[
\eta_\infty([s,\infty))=
\begin{cases}
\dfrac{1}{\sqrt5}+\dfrac{\varphi}{\sqrt5},& 0<s<\varphi^{-1},\\[2mm]
\dfrac{\varphi}{\sqrt5},& \varphi^{-1}<s<1,\\[2mm]
0,& s>1,
\end{cases}
\]
两者一致，故第一条分布恒等式成立。

再对上式取分布导数。由于 \(\widehat C_\infty'\) 只在
\[
s=\varphi^{-1},\qquad s=1
\]
发生跃迁，且两处跳跃大小分别为
\[
\frac{1}{\sqrt5},
\qquad
\frac{\varphi}{\sqrt5},
\]
故
\[
-\frac{\dd^2}{\dd s^2}\widehat C_\infty
=
\frac{1}{\sqrt5}\delta_{\varphi^{-1}}
+
\frac{\varphi}{\sqrt5}\delta_1
=
\eta_\infty.
\]

唯一性亦随之而来：若另一有限正测度 \(\eta\) 满足同一积分表示，则对函数
\[
s\longmapsto \int_{(0,\infty)}\min\{s,y\}\,\eta(\dd y)
\]
作二阶分布导数便得到
\[
\eta=-\widehat C_\infty''=\eta_\infty.
\]
故 \(\eta_\infty\) 唯一。证毕。
\end{proof}

\begin{theorem}[bin-fold 容量极限是稳定层均匀二原子律的尺寸偏置平均]\label{thm:xi-binfold-capacity-size-bias-uniform-law}
\leanverified{paper\_xi\_binfold\_capacity\_size\_bias\_uniform\_law}
沿用定理 \ref{thm:xi-foldbin-uniform-layer-degeneracy-two-atom-law} 的记号，记 \(Y_\infty^{\mathrm{unif}}\) 为该定理中的极限随机变量，即
\[
\PP\!\bigl(Y_\infty^{\mathrm{unif}}=1\bigr)=\varphi^{-1},
\qquad
\PP\!\bigl(Y_\infty^{\mathrm{unif}}=\varphi^{-1}\bigr)=\varphi^{-2}.
\]
则前一定理中的原子测度满足
\[
\boxed{\;
\eta_\infty
=
\frac{\varphi^2}{\sqrt5}\,
\mathcal L\!\bigl(Y_\infty^{\mathrm{unif}}\bigr).\;}
\]
等价地，对每个 \(s>0\) 都有
\[
\boxed{\;
\widehat C_\infty(s)
=
\frac{\varphi^2}{\sqrt5}\,
\EE\!\bigl[\min\{Y_\infty^{\mathrm{unif}},s\}\bigr].\;}
\]

进一步，记 \(Y_\infty^\mu\) 为推论 \ref{cor:fold-bin-two-point-limit-law} 的 bin-fold 推前极限随机变量。则 \(Y_\infty^\mu\) 正是 \(Y_\infty^{\mathrm{unif}}\) 的尺寸偏置变换，即
\[
\boxed{\;
\PP\!\bigl(Y_\infty^\mu=y\bigr)
=
\frac{
y\,\PP\!\bigl(Y_\infty^{\mathrm{unif}}=y\bigr)
}{
\EE\!\bigl[Y_\infty^{\mathrm{unif}}\bigr]
}
\qquad
\bigl(y\in\{1,\varphi^{-1}\}\bigr).\;}
\]
特别地，
\[
\PP\!\bigl(Y_\infty^\mu=1\bigr)=\frac{\varphi}{\sqrt5},
\qquad
\PP\!\bigl(Y_\infty^\mu=\varphi^{-1}\bigr)=\frac{1}{\varphi\sqrt5}.
\]
\end{theorem}
\begin{proof}
由定理 \ref{thm:xi-foldbin-uniform-layer-degeneracy-two-atom-law}，
\[
\mathcal L\!\bigl(Y_\infty^{\mathrm{unif}}\bigr)
=
\frac{1}{\varphi}\delta_1+\frac{1}{\varphi^2}\delta_{\varphi^{-1}}.
\]
两边同乘 \(\varphi^2/\sqrt5\)，便得到
\[
\frac{\varphi^2}{\sqrt5}\,
\mathcal L\!\bigl(Y_\infty^{\mathrm{unif}}\bigr)
=
\frac{\varphi}{\sqrt5}\delta_1+\frac{1}{\sqrt5}\delta_{\varphi^{-1}}
=
\eta_\infty,
\]
其中最后一步正是定理 \ref{thm:xi-binfold-capacity-atomic-measure-unique-recovery} 中的原子测度公式。于是
\[
\widehat C_\infty(s)
=
\int_{(0,\infty)}\min\{s,y\}\,\eta_\infty(\dd y)
=
\frac{\varphi^2}{\sqrt5}
\EE\!\bigl[\min\{Y_\infty^{\mathrm{unif}},s\}\bigr].
\]

再由定理 \ref{thm:xi-foldbin-uniform-layer-degeneracy-two-atom-law} 在 \(a=1\) 时的矩公式，
\[
\EE\!\bigl[Y_\infty^{\mathrm{unif}}\bigr]
=
\frac{1}{\varphi}+\varphi^{-3}
=
\frac{\sqrt5}{\varphi^2},
\]
其中最后一步由 \(\varphi^2=\varphi+1\) 直接化简可得。故尺寸偏置概率为
\[
\frac{1\cdot \varphi^{-1}}{\sqrt5/\varphi^2}=\frac{\varphi}{\sqrt5},
\qquad
\frac{\varphi^{-1}\cdot \varphi^{-2}}{\sqrt5/\varphi^2}
=
\frac{1}{\varphi\sqrt5}.
\]
这与推论 \ref{cor:fold-bin-two-point-limit-law} 的两点极限概率完全一致，故 \(Y_\infty^\mu\) 正是 \(Y_\infty^{\mathrm{unif}}\) 的尺寸偏置变换。证毕。
\end{proof}

\begin{theorem}[bin-fold 容量极限是统一层矩谱的 Hardy--Littlewood 变换并满足 Mellin--Stieltjes 完全对偶]\label{thm:xi-binfold-capacity-hl-mellin-duality-uniform-law}
沿用定理 \ref{thm:xi-binfold-capacity-size-bias-uniform-law} 与定理 \ref{thm:xi-foldbin-uniform-layer-two-point-exponential-family} 的记号。记
\[
Y_\infty^{\mathrm{unif}}
\sim
\frac{1}{\varphi}\delta_1+\frac{1}{\varphi^2}\delta_{\varphi^{-1}},
\qquad
Z(q):=\EE\!\left[\bigl(Y_\infty^{\mathrm{unif}}\bigr)^q\right]
=
\frac{1}{\varphi}+\varphi^{-q-2}.
\]
则：
\begin{enumerate}
  \item 对每个 \(s>0\)，bin-fold 极限容量曲线满足 Hardy--Littlewood 变换公式
  \[
  \boxed{\;
  \widehat C_\infty(s)
  =
  \frac{\varphi^2}{\sqrt5}\,
  \EE\!\bigl[\min\{Y_\infty^{\mathrm{unif}},s\}\bigr].\;}
  \]
  \item 对任意 \(0<q<1\)，有精确 Mellin--Stieltjes 对偶
  \[
  \boxed{\;
  Z(q)
  =
  \frac{\sqrt5}{\varphi^2}\,
  q(1-q)
  \int_0^\infty s^{q-2}\widehat C_\infty(s)\,\dd s.\;}
  \]
\end{enumerate}
\end{theorem}
\begin{proof}
第一条正是定理 \ref{thm:xi-binfold-capacity-size-bias-uniform-law} 的第二个盒装公式。

现证第二条。先记
\[
F(s):=\EE\!\bigl[\min\{Y_\infty^{\mathrm{unif}},s\}\bigr]
\qquad (s>0).
\]
对任意非负随机变量 \(Y\) 与任意 \(0<q<1\)，都有通用恒等式
\[
\EE[Y^q]
=
q(1-q)\int_0^\infty s^{q-2}\EE[\min\{Y,s\}]\,\dd s.
\]
事实上，由 layer-cake 公式，
\[
\EE[Y^q]
=
q\int_0^\infty s^{q-1}\PP(Y\ge s)\,\dd s,
\]
而
\[
\EE[\min\{Y,s\}]
=
\int_0^s \PP(Y\ge t)\,\dd t
\]
对 \(s\) 的导数正是 \(\PP(Y\ge s)\)。对右端作一次分部积分，即得上式。

将其应用于 \(Y=Y_\infty^{\mathrm{unif}}\)，便得到
\[
Z(q)
=
q(1-q)\int_0^\infty s^{q-2}F(s)\,\dd s.
\]
再由第一条，
\[
F(s)=\frac{\sqrt5}{\varphi^2}\,\widehat C_\infty(s),
\]
故
\[
Z(q)
=
\frac{\sqrt5}{\varphi^2}\,
q(1-q)\int_0^\infty s^{q-2}\widehat C_\infty(s)\,\dd s.
\]
证毕。
\end{proof}

\begin{theorem}[fold-gauge 缺陷序列的三次重整化正规形]\label{thm:xi-fold-gauge-defect-cubic-renormalization-normal-form}
沿用定理 \ref{thm:xi-foldbin-gauge-constant-not-cfinite-first-two-zeta-limits} 的记号，令
\[
E_m:=\log C_m-\log(2\pi),
\qquad
q:=\frac{\varphi}{2}.
\]
则当 \(m\to\infty\) 时，
\[
\boxed{\;
E_{m+1}
=
qE_m+\frac{3(2+\sqrt5)}{32}\,E_m^3+O(E_m^5).\;}
\]
等价地，
\[
\boxed{\;
\lim_{m\to\infty}
\frac{E_{m+1}-qE_m}{E_m^3}
=
\frac{3(2+\sqrt5)}{32}.\;}
\]
\end{theorem}
\begin{proof}
由定理 \ref{thm:xi-foldbin-gauge-constant-not-cfinite-first-two-zeta-limits}，
\[
E_m
=
\frac{1}{3\varphi}q^m-\frac{\sqrt5}{180}q^{3m}+O(q^{5m}).
\]
记
\[
a:=\frac{1}{3\varphi},
\qquad
b:=-\frac{\sqrt5}{180}.
\]
则
\[
E_m=aq^m+bq^{3m}+O(q^{5m}).
\]
因此
\[
E_{m+1}-qE_m
=
b(q^3-q)q^{3m}+O(q^{5m}),
\]
而
\[
E_m^3
=
a^3 q^{3m}+O(q^{5m}).
\]
于是
\[
\frac{E_{m+1}-qE_m}{E_m^3}
=
\frac{b(q^3-q)}{a^3}+O(q^{2m}).
\]
直接代入
\[
a=\frac{1}{3\varphi},
\qquad
b=-\frac{\sqrt5}{180},
\qquad
q=\frac{\varphi}{2}
\]
化简可得
\[
\frac{b(q^3-q)}{a^3}
=
\frac{3(2+\sqrt5)}{32}.
\]
故所示极限成立。

最后，由 \(a>0\) 可知
\[
E_m\sim aq^m,
\]
从而 \(q^{5m}=O(E_m^5)\)。把上式改写即得
\[
E_{m+1}
=
qE_m+\frac{3(2+\sqrt5)}{32}\,E_m^3+O(E_m^5).
\]
证毕。
\end{proof}

\endinput
